\ifdefined\gkver\else\def\gkver{student}\fi
\documentclass[\gkver]{gaokaozhenti}
\begin{document}

\chapter{1984年全国}

\begin{enumerate}
\item
%% number: 1
%% paperName: 1984年全国高考第 1 题 3 分
%% typeId: 3
%% type: 填空题
%% chapter: 曲线运动
%% point: 人造地球卫星
%% method: 无
%% score: 3
%% degree: 850
%% duplicateId: 0
%% body:
我国在 1984 年 4 月 8 日成功地发射了一颗通信卫星。这颗卫星绕地球公转的角速度 $\omega_{1}$ 跟地球自转的角速度之比 $\omega_{1}:\omega_{2} =$ \tkanswer{1:1}。
%% answer: 1:1
\jdanswer{
1:1
}
%% memo:
\memoanswer{
通信卫星是同步卫星，其绕地球转动的角速度与地球自转的角速度相同，故 $\omega_{1}:\omega_{2} = 1:1$。

评分标准：本题 3 分。
}

\item
%% number: 2
%% paperName: 1984年全国高考第 2 题 3 分
%% typeId: 3
%% type: 填空题
%% chapter: 原子和原子核
%% point: 人工核反应
%% method: 无
%% score: 3
%% degree: 850
%% duplicateId: 0
%% body:
平衡下列核反应方程式：

$\ce{^{7}_{3}Li + ^{1}_{1}H -> ^{4}_{2}He} +$ \tkanswer{$\ce{^{4}_{2}He}$}，$\ce{^{10}_{5}B} +$ \tkanswer{$\ce{^{1}_{0}n}$} $\ce{-> ^{7}_{3}Li + ^{4}_{2}He}$。
%% answer: \ce{^{4}_{2}He}，\ce{^{1}_{0}n}
\jdanswer{
$\ce{^{4}_{2}He}$，$\ce{^{1}_{0}n}$
}
%% memo:
\memoanswer{
核反应方程满足质量数与核电荷数守恒，分析可知两空分别为 $\ce{^{4}_{2}He}$、$\ce{^{1}_{0}n}$。

评分标准：本题 3 分，答对一个给 1 分，答对两个给 3 分。
}

\item
%% number: 3
%% paperName: 1984年全国高考第 3 题 3 分
%% typeId: 3
%% type: 填空题
%% chapter: 磁场
%% point: 带电粒子在复合场中的运动
%% method: 无
%% score: 3
%% degree: 750
%% duplicateId: 0
%% body:
如右图所示，一正离子以速度 $v$ 从左向右射入匀强电场和匀强磁场并存的区域中。电场强度 $E = 4\times10^{4}\UNC$，磁感应强度 $B = 0.2\UT$，方向垂直纸面向里，电场、磁场和速度三者的方向互相垂直。如果该离子在场中运动时不发生偏转，则电场方向在附图中为从\tkanswer{上}向\tkanswer{下}；离子速度大小 $v =$ \tkanswer{$2\times10^{5}\Ums$}。
\onepicture[align=right, w=5cm]{figs/03.png}
%% answer: 上，下；$2\times10^{5}\Ums$
\jdanswer{
上，下；$2\times10^{5}\Ums$
}
%% memo:
\memoanswer{
由左手定则知洛伦兹力向上，由平衡条件得电场力应向下，所以电场方向向下；由 $qvB = qE$ 得 $v = \frac{E}{B} = 2\times10^{5}\Ums$。

评分标准：本题 3 分，第一空 1 分，第二空 2 分。
}

\item
%% number: 4
%% paperName: 1984年全国高考第 4 题 3 分
%% typeId: 3
%% type: 填空题
%% chapter: 光学
%% point: 光电效应
%% method: 无
%% score: 3
%% degree: 650
%% duplicateId: 0
%% body:
频率为 $\nu$ 的光照射到一金属表面上，有电子从金属表面逸出。当所加反向电压 $U$ 的大小增大到 $3U$ 时，光电流刚好减小到零。已知这种金属的极限频率为 $\nu_{0} = 6\times10^{14}\UHz$，因此照射光的频率 $\nu =$ \tkanswer{$13.2\times10^{14}\UHz$}。
%% answer: $13.2\times10^{14}\UHz$
\jdanswer{
$13.2\times10^{14}\UHz$
}
%% memo:
\memoanswer{
由光电效应方程得 $\frac{1}{2}mv^{2} = h\nu - h\nu_{0}$，由动能定理得 $eU_{\text{止}} = \frac{1}{2}mv^{2}$，联立解得 $\nu = 13.2\times10^{14}\UHz$。

评分标准：本题 3 分。
}

\item
%% number: 5
%% paperName: 1984年全国高考第 5 题 3 分
%% typeId: 1
%% type: 单选题
%% chapter: 光学
%% point: 光谱
%% method: 无
%% score: 3
%% degree: 850
%% duplicateId: 0
%% body:
太阳光谱中含有许多暗线，这些暗线是由于\tkanswer{太阳光经过温度较低的太阳大气层时某些特征谱线的光被吸收}而形成的。
%% answer:
\jdanswer{
太阳光经过温度较低的太阳大气层时某些特征谱线的光被吸收
}
%% memo:
\memoanswer{
太阳光经过温度较低的太阳大气层时，某些特征谱线的光被吸收，在太阳光谱中形成暗线。

评分标准：本题 3 分，要求答出两个要点：太阳大气层、吸收；答出一个要点给 1 分，两个都答出给 3 分。
}

\item
%% number: 6
%% paperName: 1984年全国高考第 6 题 3 分
%% typeId: 3
%% type: 填空题
%% chapter: 原子和原子核
%% point: 玻尔模型
%% method: 无
%% score: 3
%% degree: 650
%% duplicateId: 0
%% body:
氢原子的基态能量 $E_{1} = -13.6\UeV$，则氢原子处于量子数 $n = 5$ 的能级时的能量为\tkanswer{$-0.544\UeV$}。
%% answer: $-0.544\UeV$
\jdanswer{
$-0.544\UeV$
}
%% memo:
\memoanswer{
由氢原子的能级公式 $E_{n} = \frac{E_{1}}{n^{2}}$ 得 $E_{5} = \frac{-13.6}{5^{2}}\UeV = -0.544\UeV$。

评分标准：本题 3 分。
}

\item
%% number: 7
%% paperName: 1984年全国高考第 7 题 3 分
%% typeId: 3
%% type: 填空题
%% chapter: 机械波
%% point: 干涉衍射
%% method: 无
%% score: 3
%% degree: 650
%% duplicateId: 0
%% body:
$S_{1}$ 和 $S_{2}$ 是两个相干波源。在图中分别以 $S_{1}$ 和 $S_{2}$ 为圆心作出了两组同心圆弧，分别表示在同一时刻两列波的波峰和波谷，实线表示波峰，虚线表示波谷。在图中方框内标出了三个点 a、b、c。在这三个点中，振动加强的点是\tkanswer{a、b}，振动减弱的点是\tkanswer{c}。
\onepicture[align=right, w=6cm]{figs/07.png}
%% answer: a、b；c
\jdanswer{
a、b；c
}
%% memo:
\memoanswer{
a 点为波峰与波峰相遇点，振动加强；b 点为波谷与波谷相遇点，振动加强；c 点为波峰与波谷相遇点，振动减弱。

评分标准：本题 3 分，a、b 空 2 分，c 空 1 分。
}

\item
%% number: 8
%% paperName: 1984年全国高考第 8 题 4 分
%% typeId: 1
%% type: 单选题
%% chapter: 运动和力
%% point: 牛顿第一定律
%% method: 无
%% score: 4
%% degree: 850
%% duplicateId: 0
%% body:
火车在平直轨道上匀速行驶，门窗紧闭的车厢内有一人向上跳起，发现仍落回车上原处，这是因为\xzanswer{D}
\fourchoices[answer=D]
{人跳起后，车厢内空气给他以向前的力，带着他随同火车一起向前运动}
{人跳起的瞬间，车厢地板给他一个向前的力，推动他随同火车一起向前运动}
{人跳起后，车在继续向前运动，所以人落下后必定偏后一些，只是由于时间很短，偏后距离太小，不明显而已}
{人跳起后直到落地，在水平方向上人和车始终有相同的速度}
%% answer: D
%% memo:
\memoanswer{
由牛顿第一定律知，人跳起后直到落地，在水平方向上人和车始终具有相同的速度，所以仍落回到车上原处，故 D 正确。
}

\item
%% number: 9
%% paperName: 1984年全国高考第 9 题 4 分
%% typeId: 1
%% type: 单选题
%% chapter: 电磁感应
%% point: 法拉第电磁感应定律
%% method: 无
%% score: 4
%% degree: 850
%% duplicateId: 0
%% body:
法拉第电磁感应定律可以这样表述：闭合电路中感生电动势的大小\xzanswer{C}
\fourchoices[answer=C]
{跟穿过这一闭合电路的磁通量成正比}
{跟穿过这一闭合电路的磁感应强度成正比}
{跟穿过这一闭合电路的磁通量的变化率成正比}
{跟穿过这一闭合电路的磁通量的变化量成正比}
%% answer: C
%% memo:
\memoanswer{
法拉第电磁感应定律 $E = \frac{\Delta\Phi}{\Delta t}$ 表明：闭合电路中感生电动势的大小跟穿过这一闭合电路的磁通量的变化率成正比，故 C 正确。
}

\item
%% number: 10
%% paperName: 1984年全国高考第 10 题 4 分
%% typeId: 2
%% type: 多选题
%% chapter: 电路
%% point: 动态分析
%% method: 无
%% score: 4
%% degree: 750
%% duplicateId: 0
%% body:
在下图所示的电路中，当可变电阻 $R$ 的阻值增大时，\xzanswer{AD}
\onepicture[w=6cm]{figs/10.png}
\fivechoices[answer=AD]
{AB 两点间的电压 $U$ 增大}
{AB 两点间的电压 $U$ 减小}
{通过 $R$ 的电流 $I$ 增大}
{通过 $R$ 的电流 $I$ 减小}
{通过 $R$ 的电流 $I$ 不变}
%% answer: AD
%% memo:
\memoanswer{
由闭合电路欧姆定律得，干路电流 $I = \frac{E}{R_{\text{外}} + r}$；因可变电阻 $R$ 的阻值增大，$R_{\text{外}}$ 增大，所以干路电流减小。AB 两点间的电压 $U = E - I(R_{0} + r)$，所以 $U$ 增大，故 A 正确；与可变电阻 $R$ 并联的电阻上的电流增大，可变电阻 $R$ 上的电流减小，故 D 正确。
}

\item
%% number: 11
%% paperName: 1984年全国高考第 11 题 4 分
%% typeId: 2
%% type: 多选题
%% chapter: 电场
%% point: 电场线
%% method: 无
%% score: 4
%% degree: 650
%% duplicateId: 0
%% body:
下列几种说法中，哪种说法正确？\xzanswer{AC}
\fivechoices[answer=AC]
{电场中电力线上每一点的切线方向都跟该点的场强方向一致}
{沿电力线方向，场强一定越来越小}
{沿电力线方向，电势一定越来越低}
{在电场力作用下，正电荷\textbf{一定}从电势高的地方向电势低的地方移动}
{在电场力作用下，负电荷\textbf{一定}从电势高的地方向电势低的地方移动}
%% answer: AC
%% memo:
\memoanswer{
电场中电力线上每一点的切线方向都跟该点的场强方向一致，沿电力线方向电势一定越来越低，故 A、C 正确；场强的大小由电场线的疏密决定，故 B 错误；电荷的运动方向还受其初速度影响，故 D、E 错误。
}

\item
%% number: 12
%% paperName: 1984年全国高考第 12 题 14 分
%% typeId: 4
%% type: 实验题
%% chapter: 电路
%% point: 电表改装
%% method: 无
%% score: 14
%% degree: 450
%% duplicateId: 0
%% body:
测定电流表内电阻的实验中备用的器件有：

A．电流表（量程 $0\sim100\UuA$）；

B．标准伏特表（量程 $0\sim5\UV$）；

C．电阻箱（阻值范围 $0\sim999\UO$）；

D．电阻箱（阻值范围 $0\sim9999\UO$）；

E．电源（电动势 $2\UV$，有内阻）；

F．电源（电动势 $6\UV$，有内阻）；

G．滑动变阻器（阻值范围 $0\sim50\UO$，额定电流 $1.5\UA$），还有若干电键和导线。
\begin{enumerate}
	\item
	如果采用图 1 所示的电路测定电流表 A 的内电阻并且要想得到较高的精确度，那末从以上备用的器件中，可变电阻 $R_{1}$ 应选用\tkanswer{D}，可变电阻 $R_{2}$ 应选用\tkanswer{C}，电源 $\varepsilon$ 应选用\tkanswer{F}。（用字母代号填写）
	\item
	如果实验时要进行的步骤有：

	A．合上 $K_{1}$；

	B．合上 $K_{2}$；

	C．观察 $R_{1}$ 的阻值是否最大，如果不是，将 $R_{1}$ 的阻值调至最大；

	D．调节 $R_{1}$ 的阻值，使电流表指针偏转到满刻度；

	E．调节 $R_{2}$ 的阻值，使电流表指针偏转到满刻度的一半；

	F．记下 $R_{2}$ 的阻值。

	把以上步骤的字母代号按实验的合理顺序填写在下面横线上空白处：\tkanswer{CADBEF}。
	\item
	如果在步骤 F 中所得 $R_{2}$ 的阻值为 $600\UO$，则图 1 中电流表的内电阻 $R_{g}$ 的测量值为\tkanswer{$600\UO$}。
	\item
	如果要将第（1）小题中的电流表 A 改装成量程为 $0\sim5\UV$ 的伏特表，则改装的方法是跟电流表\tkanswer{串}联一个阻值为\tkanswer{$49400\UO$}的电阻。
	\item
	图 2 所示器件中，一部分是将电流表改装为伏特表所需的，其余是为了把改装成的伏特表跟标准伏特表进行核对所需的。首先在下面空白处画出改装和核对都包括在内的电路图（要求对 $0\sim5\UV$ 的所有刻度都能在实验中进行核对），然后在图 2 上画出连线，将所示器件按以上要求连接成实验电路。
\end{enumerate}
\twopicture[labelA=1984全国12a, labelB=1984全国12b, numA=1, numB=2, labelprefix={图}, align=right, wA=6cm, wB=8cm]
{figs/12a.png}
{figs/12b.png}
%% answer: （1）D，C，F；（2）CADBEF；（3）$600\UO$；（4）串，$49400\UO$；（5）电路图与连线如图所示
\jdanswer{
\begin{enumerate}
	\item
	D，C，F

	\item
	CADBEF

	\item
	$600\UO$

	\item
	串，$49400\UO$

	\item
	电路图与连线分别如图 1、图 2 所示。
\end{enumerate}
}
%% memo:
\memoanswer{
因为电流表半偏时总电流已经改变，所以半偏已不是原来的半偏；为使总电流改变较少，须使电阻 $R_{1}$ 尽量大，所以分别选 D、C、F。

由半偏法测电流表电阻的原理知，合理的实验步骤为 CADBEF；所测电流表的电阻为 $R_{g} = R_{2} = 600\UO$。

电流表改装为电压表串联电阻的计算结果为 $R = R_{g}\left(\frac{U}{U_{g}} - 1\right) = R_{g}\left(\frac{U}{I_{g}R_{g}} - 1\right)$，代入数据得 $R = 49400\UO$。

电路图与连线分别如图 1、图 2 所示。

\onepicture[w=5cm]{figs/12_an1.png}

\onepicture[w=8cm]{figs/12_an2.png}

评分说明：本小题共 14 分。（1）（2）各 3 分，（3）1 分，（4）2 分，（5）5 分。
}

\item
%% number: 13
%% paperName: 1984年全国高考第 13 题 7 分
%% typeId: 8
%% type: 简答题
%% chapter: 运动和力
%% point: 动量守恒定律
%% method: 无
%% score: 7
%% degree: 650
%% duplicateId: 0
%% body:
根据牛顿运动定律证明：两物体沿一直线运动，相互作用但不受外力时，它们的总动量保持不变。
%% answer: 总动量保持不变
\jdanswer{
两物体相互作用时总动量的改变量等于零，总动量保持不变。
}
%% memo:
\memoanswer{
牛顿第二定律 $F = ma$ 可以写成 $F = m\frac{\Delta v}{\Delta t} = \frac{\Delta p}{\Delta t}$，其中 $p = mv$ 为物体的动量。

当两物体相互作用而不受外力时，令 $F_{1}$ 和 $F_{2}$ 分别表示每个物体所受的力，可得 $F_{1} = \frac{\Delta p_{1}}{\Delta t}$，$F_{2} = \frac{\Delta p_{2}}{\Delta t}$。

根据牛顿第三定律得 $F_{1} = -F_{2}$，因此 $\Delta p_{1} = -\Delta p_{2}$。

令 $p$ 表示两物体的总动量，则 $\Delta p = \Delta p_{1} + \Delta p_{2} = 0$，即总动量的改变量等于零，总动量保持不变。

评分说明：本小题 7 分。列出 $F_{1} = \frac{\Delta p_{1}}{\Delta t}$、$F_{2} = \frac{\Delta p_{2}}{\Delta t}$ 的给 2 分，得出 $F_{1} = -F_{2}$ 的给 1 分，得出 $\Delta p = \Delta p_{1} + \Delta p_{2} = 0$ 的再给 4 分。证明过程有正负号错误的扣 2 分。
}

\item
%% number: 14
%% paperName: 1984年全国高考第 14 题 7 分
%% typeId: 6
%% type: 计算题
%% chapter: 交流电
%% point: 变压器
%% method: 无
%% score: 7
%% degree: 750
%% duplicateId: 0
%% body:
在右图所示的电路中，一理想变压器的原线圈跟副线圈的匝数比为 $N_{1}:N_{2} = 1:2$。电源电压 $U = 220\UV$，A 是额定电流为 $I_{0} = 1\UA$ 的保险丝，$R$ 是可变电阻。为了不使原线圈中的电流超过 $I_{0}$，调节电阻 $R$ 时，其阻值最低不能小于多少欧姆？
\onepicture[align=right, w=6cm]{figs/14.png}
%% answer: $R = 880\UO$
\jdanswer{
$R = 880\UO$
}
%% memo:
\memoanswer{
设所求的电阻值为 $R$，此时 $R$ 两端的电压为 $U_{2}$（副线圈两端的电压），电流为 $I_{2}$（副线圈中的电流），原线圈中的电流为 $I_{0}$，原线圈两端的电压为 $U$，由理想变压器原理得
$UI_{0} = U_{2}I_{2}$，$U_{2} = I_{2}R$，$\frac{U_{2}}{U} = \frac{N_{2}}{N_{1}}$，
由以上各式可得 $R = \left(\frac{N_{2}}{N_{1}}\right)^{2}\frac{U}{I_{0}}$，
代入数据得 $R = 880\UO$。

评分说明：本小题 7 分。三式列对一个给 1 分，列对两个给 3 分，三个都列对给 5 分；得出 $R$ 的表达式再给 1 分，最后答数正确再给 1 分。
}

\item
%% number: 15
%% paperName: 1984年全国高考第 15 题 7 分
%% typeId: 6
%% type: 计算题
%% chapter: 力物体的平衡
%% point: 力矩平衡
%% method: 无
%% score: 7
%% degree: 550
%% duplicateId: 0
%% body:
如图为天平的原理示意图，天平横梁的两端和中央各有一刀口，图中分别用 A、B、O 三点代表；三点在一条直线上，并且 OA $=$ OB $= L$。横梁（包括固定在横梁上的指针 OD）可以中央刀口为轴转动，两边的挂架及盘的质量相等，横梁的质量为 $M$。当横梁水平时，其重心 C 在刀口的正下方，C 到 O 的距离为 $h$，此时指针竖直向下。设只在一盘中加一质量为 $\Delta m$ 的微小砝码，最后横梁在某一倾斜位置上达到平衡，此时指针与竖直方向成 $\theta$ 角。已知 $L$、$h$、$M$ 及 $\Delta m$，求 $\theta$。
\onepicture[align=right, w=5cm]{figs/15.png}
%% answer: $\tan\theta = \frac{\Delta mL}{Mh}$
\jdanswer{
$\tan\theta = \frac{\Delta mL}{Mh}$
}
%% memo:
\memoanswer{
\onepicture[align=right, w=5cm]{figs/15_an.png}

有固定转动轴物体的平衡条件是力矩的代数和等于零。设两边挂盘（包括质量为 $\Delta m$ 的微小砝码）对横梁的作用力分别为 $F_{1}$ 和 $F_{2}$，则
$F_{1}L\cos\theta = F_{2}L\cos\theta + Mgh\sin\theta$，
因为 $F_{1} - F_{2} = \Delta mg$，
所以得 $\Delta mgL\cos\theta = Mgh\sin\theta$，
即 $\tan\theta = \frac{\Delta mL}{Mh}$。

评分说明：本小题 7 分。列出 $F_{1}L\cos\theta = F_{2}L\cos\theta + Mgh\sin\theta$ 和 $\Delta mgL\cos\theta = Mgh\sin\theta$ 的给 6 分，得出 $\tan\theta = \frac{\Delta mL}{Mh}$ 再给 1 分。直接列出后式的不扣分；把倾斜后 $\Delta mg$ 的力臂 $L\cos\theta$ 写作 $L$ 的也可不扣分。
}

\item
%% number: 16
%% paperName: 1984年全国高考第 16 题 5 分
%% typeId: 6
%% type: 计算题
%% chapter: 气体性质
%% point: 气体的状态参量
%% method: 近似与估算
%% score: 5
%% degree: 650
%% duplicateId: 0
%% body:
估算地球大气层空气的总重量。（最后结果取 1 位有效数字）
%% answer: $G = 5\times10^{9}\UN$
\jdanswer{
$G = 5\times10^{9}\UN$
}
%% memo:
\memoanswer{
设地球半径为 $R$，地球表面处的大气压强为 $p$，则大气的总重量 $G = 4\pi R^{2}p$。
因 $R = 6.4\times10^{8}\Um$，$p = 1.0\times10^{5}\UPa$，代入得 $G = 5\times10^{9}\UN$。

评分说明：本小题 5 分。列出 $4\pi R^{2}p$ 给 4 分，最后答数正确再给 1 分。
}

\item
%% number: 17
%% paperName: 1984年全国高考第 17 题 13 分
%% typeId: 6
%% type: 计算题
%% chapter: 功和能
%% point: 动能定理
%% method: 无
%% score: 13
%% degree: 550
%% duplicateId: 0
%% body:
一辆车通过一根跨过定滑轮的绳 PQ 提升井中质量为 $m$ 的物体，如图所示。绳的 P 端拴在车后的挂钩上，Q 端拴在物体上。设绳的总长不变，绳的质量、定滑轮的质量和尺寸、滑轮上的摩擦都忽略不计。开始时，车在 A 点，左右两侧绳都已绷紧并且是竖直的，左侧绳长为 $H$。提升时，车加速向左运动，沿水平方向从 A 经过 B 驶向 C。设 A 到 B 的距离也为 $H$，车过 B 点时的速度为$v_{B}$。求在车由 A 移到 B 的过程中，绳 Q 端的拉力对物体做的功。
\onepicture[align=right, w=7cm]{figs/17.png}
%% answer: $W = mg(\sqrt{2} - 1)H + \frac{1}{4}mv_{B}^{2}$
\jdanswer{
$W = mg(\sqrt{2} - 1)H + \frac{1}{4}mv_{B}^{2}$
}
%% memo:
\memoanswer{
设绳的 P 端到达 B 处时，左边绳与水平地面所成夹角为 $\theta$，物体从井底上升的高度为 $h$，速度为 $v$，所求的功为 $W$，则
$W - mgh = \frac{1}{2}mv^{2}$，
因绳总长不变，所以 $h = \frac{H}{\sin\theta} - H$，$v = v_{B}\cos\theta$，
将两式代入上式，得 $W = \frac{1}{2}mv_{B}^{2}\cos^{2}\theta + mgH\left(\frac{1}{\sin\theta} - 1\right)$，
又 $\theta = \frac{\pi}{4}$，可得 $W = mg(\sqrt{2} - 1)H + \frac{1}{4}mv_{B}^{2}$。

评分说明：全题 13 分。列出 $W - mgh = \frac{1}{2}mv^{2}$ 的给 3 分，列出 $h = \frac{H}{\sin\theta} - H$ 的给 3 分，列出 $v = v_{B}\cos\theta$ 的给 5 分，列出 $\theta = \frac{\pi}{4}$ 的给 1 分，最后结果正确的再给 1 分。
}

\item
%% number: 18
%% paperName: 1984年全国高考第 18 题 10 分
%% typeId: 6
%% type: 计算题
%% chapter: 气体性质
%% point: 气体连接体
%% method: 无
%% score: 10
%% degree: 550
%% duplicateId: 0
%% body:
如图所示，在两端封闭，内径均匀的直玻璃管内，有一段汞柱将两种理想气体 a 和 b 隔开。将管竖立着，达到平衡时，若温度为 $T$，气柱 a 和 b 的长度分别为 $l_{a}$ 和 $l_{b}$；若温度为 $T'$，长度分别为 $l_{a}'$ 和 $l_{b}'$。然后将管平放在水平桌面上，在平衡时，两段气柱长度分别为 $l_{a}''$ 和 $l_{b}''$。已知 $T$、$T'$、$l_{a}$、$l_{a}'$、$l_{b}$、$l_{b}'$，求 $l_{a}''/l_{b}''$。
\onepicture[align=right, w=3cm]{figs/18.png}
%% answer: $\frac{l_{a}l_{a}'(l_{b}'T - l_{b}T')}{l_{b}l_{b}'(l_{a}'T - l_{a}T')}$
\jdanswer{
$\frac{l_{a}l_{a}'(l_{b}'T - l_{b}T')}{l_{b}l_{b}'(l_{a}'T - l_{a}T')}$
}
%% memo:
\memoanswer{
对于 a 段气体，设温度为 $T$ 时的压强为 $p_{a}$、长度为 $l_{a}$，温度为 $T'$ 时的压强为 $p_{a}'$、长度为 $l_{a}'$，由理想气体状态方程得
$\frac{p_{a}l_{a}}{T} = \frac{p_{a}'l_{a}'}{T'}$。
对于 b 段气体，同理得
$\frac{p_{b}l_{b}}{T} = \frac{p_{b}'l_{b}'}{T'}$。
由压强平衡得 $p_{b} - p_{a} = p_{b}' - p_{a}'$，将以上两式代入可解得
$\frac{p_{a}}{p_{b}} = \frac{l_{a}'(l_{b}'T - l_{b}T')}{l_{b}'(l_{a}'T - l_{a}T')}$。
管平放后两段气体的压强相等，由理想气体状态方程得 $\frac{l_{a}''}{l_{b}''} = \frac{p_{a}l_{a}}{p_{b}l_{b}}$，
所以 $\frac{l_{a}''}{l_{b}''} = \frac{l_{a}l_{a}'(l_{b}'T - l_{b}T')}{l_{b}l_{b}'(l_{a}'T - l_{a}T')}$。

评分说明：全题 10 分。四式全都列对的给 4 分；部分列对但无列错的给 1 分；有列错的不给分。压强平衡式列对给 3 分，最后结果正确再给 2 分。
}

\item
%% number: 19
%% paperName: 1984年全国高考第 19 题 10 分
%% typeId: 6
%% type: 计算题
%% chapter: 电场
%% point: 带电粒子在交变电场中的运动
%% method: 无
%% score: 10
%% degree: 450
%% duplicateId: 0
%% body:
在真空中速度为 $v = 6.4\times10^{7}\Ums$ 的电子束连续地射入两平行极板之间。极板长度为 $l = 8.0\times10^{-2}\Um$，间距为 $d = 5.0\times10^{-3}\Um$。两极板不带电时，电子束将沿两极板之间的中线通过。在两极板上加一 $50\UHz$ 的交变电压 $U = U_{0}\sin\omega t$，如果所加电压的最大值 $U_{0}$ 超过某一值 $U_{c}$ 时，将开始出现以下现象：电子束有时能通过两极板；有时间断，不能通过。
\begin{enumerate}
	\item
	求 $U_{c}$ 的大小；
	\item
	求 $U_{0}$ 为何值才能使通过的时间 $\Delta t_{\text{通}}$ 跟间断的时间 $\Delta t_{\text{断}}$ 之比为 $\Delta t_{\text{通}}:\Delta t_{\text{断}} = 2:1$。
\end{enumerate}
\onepicture[align=right, w=8cm]{figs/19.png}
%% answer: （1）$U_{c} = 91\UV$；（2）$U_{0} = \frac{182}{\sqrt{3}}\UV = 105\UV$
\jdanswer{
\begin{enumerate}
	\item
	$U_{c} = 91\UV$

	\item
	$U_{0} = \frac{182}{\sqrt{3}}\UV = 105\UV$
\end{enumerate}
}
%% memo:
\memoanswer{
（1）电子通过平行极板所用的时间 $t = \frac{l}{v} \approx 10^{-9}\Us$，交变电压的周期 $T = \frac{1}{50\UHz} = 10^{-2}\Us$，可见 $t \ll T$。因此电子通过平行极板时可认为极板间的电压不变，场强可看作匀强电场。电子沿水平方向做匀速运动，沿竖直方向做匀加速运动。设电子束刚好不能通过平行极板的电压为 $U_{c}$，则
$t = \frac{l}{v}$，$\frac{d}{2} = \frac{1}{2}at^{2}$，$a = \frac{F}{m} = \frac{eU_{c}}{md}$，
由以上三式可得 $U_{c} = \frac{mv^{2}d^{2}}{el^{2}}$，代入数值得 $U_{c} = 91\UV$。

（2）因为 $\Delta t_{\text{通}} = 2\Delta t_{\text{断}}$，由正弦函数的特性知 $U_{c} = U_{0}\sin\frac{\pi}{3}$，
由此得 $U_{0} = \frac{U_{c}}{\sin\frac{\pi}{3}}$，代入数值得 $U_{0} = \frac{182}{\sqrt{3}}\UV = 105\UV$。

评分说明：全题 10 分。（1）5 分，（2）5 分。
}

\end{enumerate}

\end{document}
